With $n=20$ the whole input space has $2^{20}\approx10^6$ points, so the disagreement between $f$ and $g$ can be computed exactly, and \cref{cor:anyD} turns compression into a bound that holds for \emph{every} input distribution. \Cref{tab:boolean_n20} uses the same architecture and training protocol as before, with widths $(20,512,512,1)$. Once a network has learned its target (test error below $10^{-5}$), its compressed description agrees with it on all $2^{20}$ inputs, or on all but at most 8 of them. Appending the exception set $E$ then costs at most about 150 bits, and the distribution-free bound is close to the uniform-distribution bound: 0.0065 against 0.0061 for \texttt{parity3} at $m=4000$, 0.018 against 0.017 for \texttt{junta6}, and 0.024 against 0.008 for \texttt{fullmaj}, whose compressed descriptions disagree with $f$ on 2--6 inputs. When $f$ still errs on a small fraction of inputs (e.g.\ \texttt{maj5} at $m=250$, with 560--3700 exceptions), the uniform-distribution bound is non-vacuous (0.12) but the distribution-free one is not: listing the exceptions costs far more bits than 250 training points can pay for. Listing $E$ costs about $|E|\log_2(2^{20}e/|E|)$ bits (7{,}000--35{,}000 here), so \cref{cor:anyD} is informative only when that is small compared with $m$. This is a limitation of the exact-equivalence route, which has to describe $f$ exactly, not of distribution-free bounds as such. For one-hidden-layer networks, the dithered transfer theorem (\cref{sec:exp-t3}) gives distribution-free bounds without exact equivalence, by certifying $f$ only at the inputs where its margins allow.

\begin{table}[t]
\centering\small
\begin{tabular}{l r r r r r r}
\toprule
task & $m$ & test err.\ $f$ & bound on $f$ & bits of $g$ & any-$D$ bound & $|E|$ \\
\midrule
\texttt{fullmaj} & 4000 & 0.0000 & 0.0078 & 41 & 0.0236 & 5 \\
\addlinespace[2pt]
\texttt{junta6} & 100 & 0.3677 & 0.8149 & 22 & 1.0000 & 270461 \\
\texttt{junta6} & 4000 & 0.0000 & 0.0168 & 93 & 0.0184 & 0 \\
\addlinespace[2pt]
\texttt{maj5} & 100 & 0.0700 & 0.3870 & 51 & 1.0000 & 73405 \\
\texttt{maj5} & 250 & 0.0022 & 0.1225 & 42 & 1.0000 & 2257 \\
\texttt{maj5} & 500 & 0.0000 & 0.0621 & 42 & 0.1237 & 3 \\
\texttt{maj5} & 1000 & 0.0000 & 0.0315 & 42 & 0.0431 & 1 \\
\texttt{maj5} & 4000 & 0.0000 & 0.0080 & 42 & 0.0083 & 0 \\
\addlinespace[2pt]
\texttt{parity3} & 100 & 0.4933 & 0.9631 & 33 & 1.0000 & 300840 \\
\texttt{parity3} & 250 & 0.4507 & 0.9118 & 449 & 1.0000 & 255715 \\
\texttt{parity3} & 500 & 0.1834 & 0.2773 & 31 & 1.0000 & 190059 \\
\texttt{parity3} & 1000 & 0.0000 & 0.0242 & 31 & 0.0545 & 2 \\
\texttt{parity3} & 4000 & 0.0000 & 0.0061 & 31 & 0.0065 & 0 \\
\bottomrule
\end{tabular}
\caption{Boolean tasks, $n=20$: test error of the trained network $f$ (273,921 parameters) and the bound on it (\cref{thm:reduction}), averaged over 3 seeds; `bits' is the description length of the compressed $g$ (including the 4-bit design charge). The any-$D$ bound (\cref{cor:anyD}) holds for every input distribution; it patches $g$ on the set $E$ of inputs where $g\ne f$ (enumerated over all $2^{20}$ inputs). The parameter-counting and VC bounds for $f$ are 1.0 in every row.}\label{tab:boolean_n20}
\end{table}

